\chapter{Zusammenfassung}

Die fortschreitende Entwicklung autonomer Fahrfunktionen führt zur Integration von Hochleistungsrechnern in moderne Fahrzeuge. Diese verfügen über erhebliche Rechenkapazitäten, die jedoch, insbesondere während der Parkphase, häufig ungenutzt bleiben. Diese Arbeit untersucht das Konzept des Vehicular Fog Computing, um diese brachliegenden Ressourcen als verteilte Recheninfrastruktur für externe Applikationen nutzbar zu machen. Da etablierte Cloud-Orchestrierungslösungen wie Kubernetes für die spezifischen Restriktionen von Fahrzeugsteuergeräten und die hohe Volatilität mobiler Netzwerke oft zu schwergewichtig sind, wird eine spezialisierte, leichtgewichtige Softwareplattform entworfen und implementiert.

Die vorgestellte Architektur basiert auf einer Trennung in einen fahrzeugseitigen \enquote{Runtime Node}, implementiert in C für maximale Hardwarenähe, und einen zentralen \enquote{Runtime Master} zur Orchestrierung. Zur Gewährleistung der Sicherheit und Portierbarkeit werden externe Applikationen mittels Linux-Container-Technologie isoliert ausgeführt. Die Kommunikation im dynamischen Netzwerkverbund wird durch den \gls{DDS} realisiert, der eine dezentrale Teilnehmererkennung ermöglicht. Ein Schwerpunkt der Arbeit liegt auf der Evaluierung geeigneter Scheduling-Strategien zur Lastverteilung. Durch simulative Untersuchungen konnte gezeigt werden, dass ein heuristischer \enquote{Greedy (TOPS Prio)}-Ansatz gegenüber exakten Verfahren eine um den Faktor 400.000 schnellere Entscheidungsfindung ermöglicht und dabei dennoch 91\,\% der optimalen Ressourcenauslastung erreicht.

Die Ergebnisse der prototypischen Implementierung belegen die technische Machbarkeit: Der entwickelte Runtime Node weist mit einer Startlatenz von unter 2 ms und einer CPU-Grundlast von unter 0,1\,\% einen minimalen Overhead auf. Zudem demonstriert die Simulation die Robustheit des Systems gegenüber Knotenausfällen, indem Aufgaben verlustfrei neu verteilt werden. Die Arbeit zeigt somit auf, dass autonome Fahrzeuge effektiv als temporäre Edge-Server in Smart-City-Szenarien integriert werden können, wenngleich zukünftige Arbeiten die Effizienz der Datenübertragung und Sicherheitsmechanismen weiter optimieren müssen.

\chapter{Abstract}

The ongoing development of autonomous driving functions is leading to the integration of high-performance computers in modern vehicles. These have considerable computing capacity, but this often remains unused, especially during parking phases. This work examines the concept of vehicular fog computing in order to make these idle resources available as a distributed computing infrastructure for external applications. Since established cloud orchestration solutions such as Kubernetes are often too heavyweight for the specific restrictions of vehicle control units and the high volatility of mobile networks, a specialised, lightweight software platform is designed and implemented.

The architecture presented is based on a separation into a vehicle-side \enquote{Runtime Node}, implemented in C for maximum hardware proximity, and a central \enquote{Runtime Master} for orchestration. To ensure security and portability, external applications are executed in isolation using Linux container technology. Communication in the dynamic network is realised by the \gls{DDS}, which enables decentralised participant recognition. One focus of the work is on evaluating suitable scheduling strategies for load distribution. Simulative investigations have shown that a heuristic \enquote{Greedy (TOPS Prio)} approach enables decision-making that is 400,000 times faster than exact methods, while still achieving 91\,\% of optimal resource utilisation.

The results of the prototype implementation prove the technical feasibility: the developed runtime node has minimal overhead with a start-up latency of less than 2 ms and a basic CPU load of less than 0.1\,\%. In addition, the simulation demonstrates the robustness of the system against node failures by redistributing tasks without loss. The work thus shows that autonomous vehicles can be effectively integrated as temporary edge servers in smart city scenarios, although future work will need to further optimise the efficiency of data transmission and security mechanisms.
